\subsection{格}

定义 1. 设 V 是域 K 上的有限维向量空间。V 关于 A 的格（或简称 V 的格）定义为满足下列条件的 V 的任一子-A-模：

存在 V 的两个自由子-A-模 \(L_{1}\)、\(L_{2}\)，使得 \(L_{1} \subset M \subset L_{2}\)，且 \(\operatorname{rg}_{\mathbf{A}}(\mathbf{L}_{1}) = \operatorname{rg}_{\mathbf{K}}(\mathbf{V})\)。

\textbf{例子}\par

\begin{enumerate}
\def\labelenumi{(\arabic{enumi})}
\item
  若取 \(\mathbf{V} = \mathbf{K}\)，则 \(\mathbf{K}\) 的格恰好是 \(\neq (0)\)\(\mathbf{K}\) 的非零分式理想（\(\S 1\)，第 1 节，定义 1）。
\item
  若 \(\operatorname{rg}_{\mathbf{K}}(\mathbf{V}) = n\)，则 V 的每个自由子-A-模 L 都有一个至多含 n 个元素的基，V 中每个在 A 上自由的子集在 K 上也自由；L 成为 V 的格的充要条件，是 L 有一个含 n 个元素的基（换言之，\(\operatorname{rg}_{\mathbf{A}}(\mathbf{L}) = n\)）。
\item
  若 A 是主理想整环，则 V 的每个格 M 都是有限生成的 A-模（因为 A 是 Noether 环），且无挠，因而是自由 A-模（代数，第七章，§ 4，第 3 节，定理 2 的推论 2）。
\end{enumerate}

命题 1. 对于 V 的子-A-模 M，要成为 V 的格，充要条件是 KM = V 且 M 被包含于 V 的某个有限生成子-A-模中。

这些条件显然必要，因为与 V 具有相同秩的 V 的自由子-A-模生成 V。反之，若 KM = V，则 M 包含 V 在 K 上的一个基 \((a_{i})_{1 \leqslant i \leqslant n}\)，因而包含由  生成的自由子-A-模 \(L_{1}\)；另一方面，若 \(a_{i}\)\(M \subset M_{1}\)，其中 \(M_{1}\) 是由有限个元素 \(b_{j}\) 生成的 V 的子-A-模，且 \((e_{i})_{1 \leqslant i \leqslant n}\) 是 V 在 K 上的一个基，则存在 A 中的元素 \(s \neq 0\)，使得每个 \(b_{j}\) 都是以 A 中元素为系数的 \(s^{-1}e_{i}\) 的线性组合；若 \(L_{2}\) 是由 \(s^{-1}e_{i}\) 生成的 V 的自由子-A-模，则 \(M \subset L_{2}\)。

推论。设 A 是 Noether 环；对于 V 的子-A-模 M，要成为 V 的格，充要条件是 KM = V 且 M 是有限生成的。

注（1）。回忆，对于 V 的每个子-A-模 M，典范映射 \(M \otimes_{A} K \to V\) 是单射，且像为 KM（代数，第二章，§ 7，第 10 节，命题 26）；因此说 KM = V 就意味着此映射是双射。

命题 2. 设 M 是 V 的格，\(M_{1}\) 是 V 的子-A-模。若存在 \(K^{*}\) 中的两个元素 x、y，使得 \(xM \subset M_{1} \subset yM\)，则 \(M_{1}\) 是 V 的格；反之，若 \(M_{1}\) 是 V 的格，则存在 A 中的两个非零元 a、b，使得 \(aM \subset M_{1} \subset b^{-1}M\)。

若 \(L_{1}\)、\(L_{2}\) 是 V 的两个自由格，满足 \(L_{1} \subset M \subset L_{2}\)，则关系 \(xM \subset M_{1} \subset yM\) 推出 \(xL_{1} \subset M_{1} \subset yL_{2}\)，且 \(xL_{1}\) 和 \(yL_{2}\) 是自由格；反之，若 \(M_{1}\) 是一个格，且 \((e_{i})_{1 \leqslant i \leqslant n}\) 是 \(L_{2}\) 在 A 上的一个基，则关系 \(KM_{1} = V\) 推出存在 \(x = a/s \in K^{*}\)（其中 a、s 是 A 中的非零元），使得对所有 i 有 \(xe_{i} \in M_{1}\)，从而 \(xM \subset xL_{2} \subset M_{1}\)，特别地 \(aM \subset M_{1}\)；交换 M 与 \(M_{1}\) 的角色，同理可见存在 A 中的 \(b \neq 0\)，使得 \(bM_{1} \subset M\)。

命题 3. (i) 若 \(\mathbf{M}_1\) 和 \(\mathbf{M}_2\) 是 \(\mathbf{V}\) 的格，则 \(\mathbf{M}_1 \cap \mathbf{M}_2\) 和 \(\mathbf{M}_1 + \mathbf{M}_2\) 也是格。

\begin{enumerate}
\def\labelenumi{(\roman{enumi})}
\setcounter{enumi}{1}
\item
  若 \(\mathbf{W}\) 是 \(\mathbf{V}\) 的向量子空间，且 \(\mathbf{M}\) 是 \(\mathbf{V}\) 的格，则 \(\mathbf{M} \cap \mathbf{W}\) 是 \(\mathbf{W}\) 的格。
\item
  设 \(\mathbf{V},\mathbf{V}_1,\ldots ,\mathbf{V}_k\) 是 \(\mathbf{K}\) 上的有限秩向量空间，并令
\end{enumerate}

\[
f \colon \mathrm{V} _ {1} \times \dots \times \mathrm{V} _ {k} \rightarrow \mathrm{V}
\]

为一个像生成 V 的多重线性映射。若 \(M_{i}\) 是 \(V_{i}\) 的格（\(1 \leqslant i \leqslant k\)），则由 \(f(\mathbf{M}_{1} \times \cdots \times \mathbf{M}_{k})\) 生成的 V 的子-A-模是 V 的格。

\begin{enumerate}
\def\labelenumi{(\roman{enumi})}
\setcounter{enumi}{3}
\item
  设 \(\mathbf{V}\) 和 \(\mathbf{W}\) 是 \(\mathbf{K}\) 上的两个有限秩向量空间，\(\mathbf{M}\) 是 \(\mathbf{V}\) 的格，\(\mathbf{N}\) 是 \(\mathbf{W}\) 的格。 的子-A-模 \(\mathbf{N}:\mathbf{M}\) 由满足  的 K-线性映射  组成，是 \(\operatorname{Hom}_{\mathbf{K}}(\mathbf{V},\mathbf{W})\) 的格。\(f\)\(f(\mathbf{M})\subset \mathbf{N}\)\(\operatorname{Hom}_{\mathbf{K}}(\mathbf{V},\mathbf{W})\)
\item
  根据命题 2，存在  中的非零元 \(a\) 和 \(b\)，使得 \(\mathbf{A}\)\(a\mathbf{M}_1 \subset \mathbf{M}_2 \subset b^{-1}\mathbf{M}_1\)；于是 \(\mathbf{M}_1 \cap \mathbf{M}_2\) 和 \(\mathbf{M}_1 + \mathbf{M}_2\) 都位于 \(a\mathbf{M}_1\) 与 \(b^{-1}\mathbf{M}_1\) 之间，因而根据命题 2 它们是格。
\item
  令 S 是 V 中 W 的一个补，\(L_{w}\) 是 W 的自由格，\(L_{s}\) 是 S 的自由格，从而 \(L = L_{w} \oplus L_{s}\) 是 V 的自由格。于是存在 \(K^{*}\) 中的 x、y，使得 \(xL \subset M \subset yL\)。由此得到 \(xL_{w} \subset M \cap W \subset yL_{w}\)，这说明 \(M \cap W\) 是 W 的格（命题 2）。
\item
  由于 \(KM_{i} = V_{i}\)，由线性性显然 \(f(\mathbf{M}_{1} \times \cdots \times \mathbf{M}_{k})\) 生成 K-向量空间 V；另一方面，对所有 i，都存在  的有限生成子-A-模 \(N_{i}\)，使得 \(V_{i}\)\(M_{i} \subset N_{i}\)；由 \(f(\mathbf{N}_{1} \times \cdots \times \mathbf{N}_{k})\) 生成的 V 的子-A-模 N 是有限生成的并包含 M，因而 M 是 V 的格（命题 1）。
\item
  令 P（相应地 Q）为包含 M（相应地包含于 N）的 V（相应地 W）的自由格；显然 N: M ⊃ Q: P。现在立即可见 Q: P 同构于 Hom \(_{A}\) (P, Q)，因而是秩为 (rg \(_{A}\) P)(rg \(_{A}\) Q) 的自由 A-模（代数，第二章，§ 1，第 6 节，命题 6 的推论 1），所以是 Hom \(_{K}\) (V, W) 的格。同样，若 P'（相应地 Q'）是包含于 M（相应地包含 N）的 V（相应地 W）的自由格，则 Q': P' ⊃ N: M，且 Q': P' 是 Hom \(_{K}\) (V, W) 的格；结论得出。
\end{enumerate}

\textbf{注}\par

\begin{enumerate}
\def\labelenumi{(\arabic{enumi})}
\setcounter{enumi}{1}
\item
  命题 3 (i) 表明，V 的格集合 R(V) 关于包含关系构成格序集；此外，若 M 是 V 的一个固定格，则 xM（其中 x 遍历 K*）构成 R(V) 的一个子集，该子集既共初又共终（集合论，第三章，§ 1，第 7 节）。
\item
  在命题 3 (iv) 的记号下，典范映射
\end{enumerate}

\[
\mathrm{N}: \mathrm{M} \rightarrow \operatorname{Hom} _ {\mathrm{A}} (\mathrm{M}, \mathrm{N}),
\]

它把每个 K-线性映射 \(f \in N\) : M 映为从 M 到 N 的 A-线性映射，该映射与 \(f|_{M}\) 具有相同的图；这个映射是双射。每个 A-线性映射 \(g\) : \(M \to N\) 都可嵌入一个 K-线性映射

\[
g \otimes 1: M \otimes_ {A} K \rightarrow N \otimes_ {A} K
\]

并且已经看到，\(M \otimes_{A} K\) 和 \(N \otimes_{A} K\) 分别与 V 和 W 识别。

特别地，若取 W = K、N = A，则 \(\operatorname{Hom}_{K}(V, W)\) 就是 V 的对偶 K-向量空间 \(V^{*}\)，而 A: M 可识别为 M 的对偶 A-模 \(M^{*}\)；今后作此识别，并称 \(M^{*}\) 为 M 的对偶格：因此它是满足对所有  都有  的 \(x^{*} \in V^{*}\) 所成的集合。\(\langle x, x^{*} \rangle \in A\)\(x \in M\)

推论。设 U、V、W 是 K 上的三个有限秩向量空间，且 \(f: \mathbf{U} \times \mathbf{V} \to \mathbf{W}\) 是左非退化的 K-双线性映射（代数，第九章，§ 1，第 1 节，定义 3）。若 M 是 V 的格，N 是 W 的格，则集合 N: \(f\mathbf{M}\)，即由满足对所有  都有  的 \(x \in \mathbf{U}\) 组成的集合，是 U 的格。\(f(x, y) \in \mathbf{N}\)\(y \in \mathbf{M}\)

令 \(s_f \colon \mathbf{U} \to \operatorname{Hom}_{\mathbf{K}}(\mathbf{V}, \mathbf{W})\) 为与 \(f\) 左相伴的 K-线性映射（代数，第九章，同上），使得 \(s_f(x)\) 是映射 \(y \mapsto f(x, y)\)；回忆，\(f\) 左非退化意味着 \(s_f\) 是单射。根据命题 3 (iv)，\(\mathbf{N} \colon \mathbf{M}\) 是 \(\operatorname{Hom}_{\mathbf{K}}(\mathbf{V}, \mathbf{W})\) 的格；由于 \(\mathbf{N} \colon {}_f \mathbf{M} = s_f^{-1}(\mathbf{N} \colon \mathbf{M})\) 且 \(s_f\) 是单射，推论由命题 3 (ii) 得出。

\textbf{例子}\par

\begin{enumerate}
\def\labelenumi{(\arabic{enumi})}
\setcounter{enumi}{3}
\tightlist
\item
  设 S 是一个有限秩的（不必结合的）K-代数，带有单位元；则从 S × S 到 S 的双线性映射 \((x,y)\mapsto xy\) 是（左、右）非退化的。若 M、N 是 S 关于 A 的格，则 M.N（命题 3 (iii)）以及由满足  的 \(x\in S\) 组成的集合（命题 3 的推论）也是格。注意，S 中存在一个包含 S 的单位元且为 S 的格的子-A-代数；为此，取 S 的一个基 \(xM\subset N\)\((e_{i})_{1\leqslant i\leqslant n}\)，使 \(e_{1}\) 为 S 的单位元，并令 \(e_{i}e_{j}=\sum_{k}c_{ijk}e_{k}\) 为 S 的乘法表（\(1\leqslant i\leqslant n,1\leqslant j\leqslant n\)），从而 \(c_{1jk}=\delta_{jk},c_{11k}=\delta_{1k}\)（克罗内克符号）。令 \(s\in A\) 为非零元，并使 \(c_{ijk}^{\prime}=s.c_{ijk}\in A\) 对所有指标三元组 \((i,j,k)\) 都成立；若对  写 \(e_{i}^{\prime}=s^{-1}e_{i}\)，则\(i\geqslant2\)
\end{enumerate}

\[
e _ {i} ^ {\prime} e _ {j} ^ {\prime} = s c _ {i j 1} ^ {\prime} e _ {1} + \sum_ {k \geqslant 2} c _ {i j k} ^ {\prime} e _ {k} ^ {\prime}
\]

当 \(i \geqslant 2\) 且 \(j \geqslant 2\) 时；以 \(e_1\) 和 \(e_i' (2 \leqslant i \leqslant n)\) 为基的 S 的格是 S 的一个子-A-代数，单位元为 \(e_1\)。

\begin{enumerate}
\def\labelenumi{(\arabic{enumi})}
\setcounter{enumi}{4}
\tightlist
\item
  设 V 是 K 上的有限维向量空间，f 是 V 上的非退化双线性型。若 M 是 V 的格，则由命题 3 的推论，集合 \(M_{f}^{*}\)，即由满足对所有  都有  的 \(x \in V\) 组成的集合，也是 V 的格；若 \(f(x, y) \in A\)\(y \in M\)\(s_{f}: V \to V^{*}\) 是与 f 左相伴的线性映射（它是双射），则 \(s_{f}(M_{f}^{*})\) 正是 M 的对偶格 \(M^{*}\)。
\end{enumerate}

命题 4. 设 B 是整环，A 是 B 的子环，K 和 L 分别是 A 和 B 的分式域。设 V 是 K 上的有限维向量空间。

\begin{enumerate}
\def\labelenumi{(\roman{enumi})}
\item
  对于 \(\mathbf{M}\)\(\mathbf{V}\) 关于 \(\mathbf{A}\) 的每个格 \(\mathbf{BM}\)，\(\mathbf{M}_{(\mathbf{B})} = \mathbf{M} \otimes_{\mathbf{A}} \mathbf{B}\) 在 \(\mathbf{V}_{(\mathbf{L})} = \mathbf{V} \otimes_{\mathbf{K}} \mathbf{L}\) 中的像 ，是 \(\mathbf{V}_{(\mathbf{L})}\) 关于 \(\mathbf{B}\) 的格。
\item
  进一步设 B 是平坦 A-模。则典范映射 \(\mathbf{M}_{(\mathbf{B})}\to\mathbf{BM}\) 是双射。若 B 还是忠实平坦的，则把 V 关于 A 的每个格 M 映为 \(\mathbf{V}_{(\mathbf{L})}\) 关于 B 的格 BM 的映射是单射。
\item
  由于 \(\mathbf{KM} = \mathbf{V}\)，显然 \(\mathbf{L}\) . (BM) = \(\mathbf{V}_{(L)}\)；另一方面，\(\mathbf{M}\) 被包含于  的有限生成子-A-模 \(\mathbf{M}_1\) 中，因而 \(\mathbf{V}\)\(\mathbf{BM}\) 被包含于 \(\mathbf{BM}_1\)，后者是有限生成的 B-模；断言 (i) 由此得出（命题 1）。
\item
  \(V_{(L)} = V \otimes_{K} L = V \otimes_{A} L\)（第二章，§ 2，第 7 节，命题 18），且由于 L 是平坦 B-模，它也是平坦 A-模（第一章，§ 2，第 7 节，命题 8 的推论 3）。由于 B 是平坦 A-模，典范映射 \(M \otimes_{A} B \to V \otimes_{A} B\) 是单射；另一方面，由于 V 是自由 K-模且 K 是平坦 A-模，V 是平坦 A-模（第一章，§ 2，第 7 节，命题 8 的推论 3），从而典范映射 \(V \otimes_{A} B \to V \otimes_{A} L\) 是单射，这证明了第一断言。还要说明，当 B 是忠实平坦 A-模时，对于 V 关于 A 的两个格 ，关系 \(BM_{1} = BM_{2}\) 推出 \(M_{1} = M_{2}\)，先注意 \(M_{1}, M_{2}\)\(BM_{1} \cap BM_{2} = B(M_{1} \cap M_{2})\)（第一章，§ 2，第 6 节，命题 6）；因此可限于 \(M_{1} \subset M_{2}\) 的情形，而我们的断言由第一章 § 3 第 1 节命题 3 应用于典范单射 \(M_{1} \to M_{2}\) 得出。
\end{enumerate}

推论。设 A 是离散赋值环。令 \(\hat{\mathbf{A}}\) 为其完备化，令 \(\hat{\mathbf{K}}\) 为 \(\hat{\mathbf{A}}\) 的分式域（第六章，§ 5，第 3 节）。映射 \(\phi\) 把 V 的每个格 M 映为 \(\hat{\mathbf{A}}\mathbf{M}\)\(\hat{\mathbf{V}} = \mathbf{V} \otimes_{\mathbf{K}} \hat{\mathbf{K}}\) 关于 \(\hat{\mathbf{A}}\) 的格 ，它是双射；其逆把 \(\hat{\mathbf{V}}\) 关于 \(\hat{\mathbf{A}}\) 的每个格 M' 映为 M' \(\cap\) V，其中 V 已典范地识别为 \(\hat{\mathbf{V}}\) 的一个子-K-向量空间。

若 L 是 V 的自由格，则 aL（其中 \(a \in A\)、\(a \neq 0\)）构成 V 上某个拓扑 T 的 0 的邻域基（与其 A-模结构相容）；取 L 在 A 上的一个基后，该拓扑可识别为 \(K^{n}\) 上的乘积拓扑；根据命题 2，T 的 0 的邻域基也由 V 关于 A 的所有格组成；显然，\(\hat{V}\) 是 V 关于 T 的完备化。此外，若 m 是 A 的极大理想，则 T 在 V 关于 A 的每个格 M 上诱导出 m-进拓扑，因为 M 是有限生成的 A-模（第三章，§3，第 2 节，定理 2），且 \(\hat{A}M\) 是 M 关于此拓扑的完备化（第三章，§2，第 12 节，命题 16）；又由于 M 在 V 中是开集（因而是闭集），有 \(\hat{A}M \cap V = M\)，这再次证明 \(\phi\) 是单射（也可直接由命题 4 (ii) 得到，因为 \(\hat{A}\) 是忠实平坦 A-模）。最后，若 \(M'\) 是 \(\hat{V}\) 关于 \(\hat{A}\) 的格，则 \(M = M' \cap V\) 是 V 关于 A 的格，因为 \(\hat{A}\) 的每个元素都是 A 中一个元素与 \(\hat{A}\) 中一个可逆元的乘积，因而由命题

2 可知存在 \(a, b\) 属于 \(\mathbf{A} - \{0\}\)，使得 \(a\hat{\mathbf{A}}\mathbf{L} \subset \mathbf{M}' \subset b\hat{\mathbf{A}}\mathbf{L}\)，从而 \(a\mathbf{L} \subset \mathbf{M}' \cap \mathbf{V} \subset b\mathbf{L}\)。此外，\(\mathbf{M}'\) 在 \(\mathbf{V}\) 中是开集，且由于 \(\mathbf{V}\) 在 \(\hat{\mathbf{V}}\) 中稠密，\(\mathbf{M}'\) 是 \(\mathbf{M}' \cap \mathbf{V} = \mathbf{M}\) 的完备化；这证明 \(\phi\) 是满射，推论得证。

例（6）。设 S 是 A 中不含 0 的乘性子集；将命题 4 应用于 \(B = S^{-1}A\)；此时 L = K，\(BM = S^{-1}M\)；因此 \(S^{-1}M\) 是 V 关于 \(S^{-1}A\) 的格。此外：

命题 5. 设 V、W 是 K 上的有限秩向量空间，M 是 V 的格，N 是 W 的格。若 M 是有限生成的，则（沿用命题 3 的记号）：

\[
\mathbf {S} ^ {- 1} (\mathbf {N}: \mathbf {M}) = \mathbf {S} ^ {- 1} \mathbf {N}: \mathbf {S} ^ {- 1} \mathbf {M}\tag{1}
\]

在 \(\operatorname{Hom}_{\mathbf{K}}(\mathbf{V},\mathbf{W})\) 中

显然，(1) 的左边包含于右边。反之，令 \(f \in S^{-1}N: S^{-1}M\)，并令 \((x_i)_{1 \leqslant i \leqslant n}\) 为 M 的一个生成元系。存在 \(s \in S\)，使得对所有  都有 \(f(x_i) \in s^{-1}N\)，因而 \(i\)\(sf \in N: M\)，命题得证。

